%! Characteristics of Logarithmic Functions — Unit 8, Lesson 8.0 — Homework (Answer Key)
\documentclass[10pt]{article}
\usepackage{algebra2-article}
\usepackage{algebra2-key}   % pulls in -boxes; do NOT also load -boxes
\newcommand{\TallMath}[1]{$\displaystyle #1\rule[-1.4em]{0pt}{3.2em}$}
\newcommand{\lgrid}[4]{%
  \draw[step=1, linegray!40, very thin] (#1,#3) grid (#2,#4);
  \draw[->, semithick, charcoal] (#1,0)--(#2,0) node[below right, font=\scriptsize]{$x$};
  \draw[->, semithick, charcoal] (0,#3)--(0,#4) node[left, font=\scriptsize]{$y$};}
% pgfmath has no log_b, so every logarithmic curve is drawn by parameterizing on y:
%   y = log_b(x-h) + k   <=>   x = b^(y-k) + h.   Args: {b}{h}{k}{ymin}{ymax}
\newcommand{\logcurve}[5]{\draw[very thick, forest, domain=#4:#5, samples=120, smooth]
  plot ({pow(#1,\x-(#3))+(#2)},{\x});}

\begin{document}

\pageheader{Unit 8, Lesson 8.0}{Homework --- Answer Key}
\namedateperiod

\vspace{0.10in}

\begin{notesbox}{Practice}
Show how you read each feature. Write every domain and range in \textbf{interval notation}, and
remember that a logarithm's boundary is never included.

\begin{enumerate}[leftmargin=1.9em, itemsep=12pt]
  \item \textbf{From the equation only.} Use the argument --- no graphing.
    \begin{center}
    \renewcommand{\arraystretch}{1.6}
    \begin{tabularx}{\linewidth}{@{}l X X X X@{}}
    \hline
    \textbf{Function} & \textbf{Condition on the argument} & \textbf{Domain} & \textbf{Vertical asymptote} & \textbf{Runs left or right?} \\
    \hline
    $\log_3(x-2)$ & \ans{$x-2>0$} & \ans{$(2,\infty)$} & \ans{$x=2$} & \ans{right} \\
    $\log(x+7)$ & \ans{$x+7>0$} & \ans{$(-7,\infty)$} & \ans{$x=-7$} & \ans{right} \\
    $\log_2(4x)$ & \ans{$4x>0$} & \ans{$(0,\infty)$} & \ans{$x=0$} & \ans{right} \\
    $\log_5(10-2x)$ & \ans{$10-2x>0$} & \ans{$(-\infty,5)$} & \ans{$x=5$} & \ans{left} \\
    \hline
    \end{tabularx}
    \end{center}

  \item \textbf{A shift that moves the intercept but not the asymptote.} The graph is
        $m(x) = \log_2 x - 3$.
    \begin{center}
    \begin{tikzpicture}[scale=0.44]
      \lgrid{-2}{13}{-6}{2}
      \draw[navy, dashed, semithick] (0,-6)--(0,2);
      \begin{scope}
        \clip (-2,-6) rectangle (13,2);
        \logcurve{2}{0}{-3}{-6}{0.70}
      \end{scope}
      \fill[charcoal] (1,-3) circle (3pt) node[below right, font=\scriptsize]{$(1,-3)$};
      \fill[charcoal] (2,-2) circle (3pt);
      \fill[charcoal] (4,-1) circle (3pt) node[above left, font=\scriptsize]{$(4,-1)$};
      \fill[charcoal] (8,0) circle (3pt) node[above left, font=\scriptsize]{$(8,0)$};
    \end{tikzpicture}
    \end{center}
    Domain: \ans{$(0,\infty)$} \quad Range: \ans{all reals} \quad Vertical asymptote:
    \ans{$x=0$} \\[3pt]
    $x$-intercept: \ans{$(8,0)$} \quad $y$-intercept: \ans{none} \\[3pt]
    Increasing or decreasing on its domain? \ans{increasing} \quad Absolute extrema: \ans{none} \\[3pt]
    $m(x)<0$ on \ans{$(0,8)$}; \quad $m(x)>0$ on \ans{$(8,\infty)$} \\[3pt]
    End behavior: as $x\to+\infty$, $m(x)\to$ \ans{$+\infty$}. \; As $x\to$ \ans{$0^{+}$},
    $m(x)\to$ \ans{$-\infty$} \\[3pt]
    The parent $y=\log_2 x$ has $x$-intercept $(1,0)$ and asymptote $x=0$. Subtracting $3$ moved
    \emph{one} of those two. Which one, and why not the other? \ans{Only the $x$-intercept moved
    (from $(1,0)$ to $(8,0)$). Subtracting 3 changes the \emph{outputs}, and the asymptote is
    determined by which \emph{inputs} are legal --- the argument is still just $x$, so the domain and
    the asymptote are unchanged.}

  \item \textbf{A base between 0 and 1.} The graph is $n(x) = \log_{1/3} x$.
    \begin{center}
    \begin{tikzpicture}[scale=0.44]
      \lgrid{-2}{11}{-3}{3}
      \draw[navy, dashed, semithick] (0,-3)--(0,3);
      \begin{scope}
        \clip (-2,-3) rectangle (11,3);
        \logcurve{0.333333}{0}{0}{-2.19}{3}
      \end{scope}
      \fill[charcoal] (1,0) circle (3pt) node[above right, font=\scriptsize]{$(1,0)$};
      \fill[charcoal] (3,-1) circle (3pt) node[below right, font=\scriptsize]{$(3,-1)$};
      \fill[charcoal] (9,-2) circle (3pt) node[below right, font=\scriptsize]{$(9,-2)$};
    \end{tikzpicture}
    \end{center}
    Domain: \ans{$(0,\infty)$} \quad Range: \ans{all reals} \quad Vertical asymptote:
    \ans{$x=0$} \\[3pt]
    $x$-intercept: \ans{$(1,0)$} \quad $y$-intercept: \ans{none} \\[3pt]
    Increasing or decreasing? \ans{decreasing} \quad Absolute extrema: \ans{none} \\[3pt]
    End behavior: as $x\to+\infty$, $n(x)\to$ \ans{$-\infty$}; \; as $x\to$ \ans{$0^{+}$},
    $n(x)\to$ \ans{$+\infty$} \\[3pt]
    Compared with $y=\log_3 x$ from class, list every feature that is the \textbf{same}:
    \ans{domain, range, vertical asymptote, $x$-intercept, no $y$-intercept, no extrema --- only the
    direction (and therefore the end behavior and the positive/negative intervals) differs}

  \item \textbf{Explain.} Every parent logarithmic graph $y=\log_b x$ passes through $(1,0)$ and
        never crosses the $y$-axis --- no matter what the base is. Explain both facts, using the
        words \textbf{exponent} and \textbf{asymptote}.
        \ansline{$\log_b 1$ asks ``$b$ to what \textbf{exponent} gives $1$?'' and the answer is
        always $0$, because $b^0=1$ for every legal base. So $(1,0)$ is on every parent log
        graph.\\[3pt]
        Crossing the $y$-axis would mean evaluating the function at $x=0$, which asks ``$b$ to what
        \textbf{exponent} gives $0$?'' No exponent does --- powers of a positive base are always
        positive --- so $x=0$ is outside the domain. The $y$-axis is the graph's vertical
        \textbf{asymptote}: the curve dives alongside it forever and never reaches it.}
\end{enumerate}
\end{notesbox}

\begin{extensionbox}
\textbf{Building a log graph from an exponential table.} Below is a table for $y=3^x$. You will not
need a calculator.
\begin{center}
\renewcommand{\arraystretch}{1.25}
\begin{tabular}{|c|c|c|c|c|c|}
\hline
$x$ & $-2$ & $-1$ & $0$ & $1$ & $2$ \\ \hline
$y=3^x$ & $\tfrac19$ & $\tfrac13$ & $1$ & $3$ & $9$ \\ \hline
\end{tabular}
\end{center}
\begin{enumerate}[label=(\alph*), leftmargin=1.8em, itemsep=5pt]
  \item Trade the rows to build a table for $y=\log_3 x$. Which row becomes the inputs?
  \item Using only your new table, state the domain and range of $y=\log_3 x$. Explain how each one
        came from a feature of $y=3^x$.
  \item The graph of $y=3^x$ has a horizontal asymptote at $y=0$. Reflect that line over $y=x$: what
        line do you get, and what is it called on the logarithmic graph?
  \item $y=3^x$ has a $y$-intercept at $(0,1)$ and no $x$-intercept at all. Use the coordinate swap
        to predict both intercept facts for $y=\log_3 x$ --- then check them against your table.
  \item In Unit~6, $y=x^2$ had to be cut down to $x\ge 0$ before $\sqrt{x}$ could undo it, but
        $y=3^x$ needs no such cut. What is true about an exponential graph that makes the difference?
\end{enumerate}
\ansline{\textbf{(a)} The old \emph{output} row becomes the inputs:
$x = \tfrac19,\ \tfrac13,\ 1,\ 3,\ 9$ with $y = -2,\ -1,\ 0,\ 1,\ 2$.}
\ansline{\textbf{(b)} Domain $(0,\infty)$ --- every new input came from an output of $3^x$, and those
are all positive. Range: all reals --- the new outputs are the old inputs, and any real number is a
legal exponent.}
\ansline{\textbf{(c)} Reflecting $y=0$ over $y=x$ gives $x=0$: the \emph{vertical asymptote} of
$y=\log_3 x$. The same boundary, turned a quarter turn.}
\ansline{\textbf{(d)} Swapping $(0,1)$ gives $(1,0)$ --- the $x$-intercept, which the table confirms.
``No $x$-intercept'' swaps into ``no $y$-intercept,'' and indeed $x=0$ never appears in the new input
row.}
\ansline{\textbf{(e)} An exponential graph is always increasing (or always decreasing), so it never
sends two different inputs to the same output. A parabola does ($(-3)^2=3^2$), which is why half of
it had to be discarded before it could be reflected into a function. All of $y=3^x$ can be reflected
as it stands.}
\end{extensionbox}

\begin{spiralbox}
\textbf{Coming next --- Lesson 8.1: Introduction to Logarithms.} Today $\log_b x$ was the
\emph{name of a curve}: you read it off a graph and used it to decide which inputs were legal. Next
you will learn to \emph{compute} with it --- writing $\log_b x = y$ and $b^y = x$ as one statement
read two ways, evaluating logs by asking ``$b$ to what power?'', meeting the common log, and using
change of base so a calculator can handle any base at all. That is also where today's domain rule
gets its full algebraic argument.
\end{spiralbox}

\begin{teachernote}
\textbf{Item 2 is the one to grade closely.} It separates students who understand that the argument
controls the domain from students pattern-matching on ``there is a number in the equation, so
something moves.'' A vertical shift changes outputs only; the asymptote does not budge. Lesson~8.2
formalizes this ($f(x+k)$ is the only transformation that moves a vertical asymptote --- the exact
mirror of what $f(x)+k$ did to the horizontal asymptote in 7.2), so a student who reasons it out here
arrives at 8.2 with the idea already built.

\textbf{Common wrong answers.} Item~1 row~4: $(-\infty,5)$ is missed by students who forget the flip
when dividing by $-2$, or who write $(-\infty,5]$. Item~3: students frequently write the domain as
``all reals'' because the graph looks like it covers everything --- point them back at the picture on
the \emph{left} of the asymptote, where nothing is drawn.

\textbf{The Extension is the lesson's $\bullet$ spine row done independently.} If a student can do
parts (c)--(e) unaided, they have the inverse relationship. Part (e) is the Unit~6 callback and is
worth a full-credit conversation even when phrased loosely (``it never repeats a $y$-value'').
\end{teachernote}

\end{document}
